\شروع‌مسئله[17]
یک مجموعه از اعداد را $k$-محسوس می‌گوییم اگر جمع هر زیرمجموعۀ $k$ عضوی از آن بزرگ‌تر یا مساوی هر کدام از اعضای مجموعه باشد. فرض کنید $M$ یک مجموعه از اعداد متمایز است و می‌خواهیم بزرگترین زیرمجموعۀ $k$-محسوس آن را پیدا کنیم. فرض کنید بزرگ‌ترین زیرمجموعۀ $k$-محسوس از $M$ یکتا است. در پاسخ به موارد زیر اثبات دقیق ادعای خود را بیان کنید.
\شروع{زیرمسائل}
\زیرمسئله[7]
ثابت کنید بزرگ‌ترین زیرمجموعۀ $k$-محسوس از $M$ محدب است، یعنی اگر شامل $a$ و $b$ باشد، شامل هر عضو $c$ که $a<c<b$ و $c\in M$ هم باشد.

\پاسخ{
	فرض کنید $A$ بزرگ‌ترین زیرمجموعۀ $k$-محسوس از $M$ و $s$ مجموع $k$ عضو کوچکتر $A$ باشد. طبق فرض $k$-محسوس بودن $A$، $s$ بزرگ‌تر یا مساوی هر کدام از عناصر $A$ است. همچنین بنابر تعریف $s$، مجموع هر $k$ عضو $A$ حداقل $s$ است. بنابراین شرط لازم و کافی $k$-محسوس بودن $A$ این است که اعضای آن کوچکتر از $s$ نباشند. از طرفی اعضای $s$ در $M$ متوالی هستند، در غیر این صورت می‌توان کوچکترین عنصر $A$ را حذف و یکی از عناصر جا افتاده را اضافه کرد و $A$ جدید همچنان $k$-محسوس است. این موضوع با یکتا بودن $A$ در تناقض است.
	
حال فرض خلف کنید که $A$ محدب نیست، بنابراین سه عنصر $a<b<c$ از $M$ را درنطر بگیرید که $a,c\in A$ و $b\notin A$. از طرفی $s\leq a<b$ پس $b$ را می‌توان به $A$ اضافه کرد و شرط $k$-محسوس بودن باقی می‌ماند (طبق متوالی بودن اعضای $s$، اضافه کردن عنصر $b$ روی $s$ تأثیر نمی‌گذارد). بنابراین، $A$ بزرگ‌ترین زیرمجموعۀ $k$-محسوس از $M$ نیست و فرض خلف غلط است. پس $A$ محدب است.
	
}
\زیرمسئله[10]
الگوریتمی با زمان اجرای $O(n\log n)$ طراحی کنید که بزرگ‌ترین زیرمجموعۀ $k$-محسوس از $M$ را پیدا کند. درستی و زمان اجرای الگوریتم خود را ثابت کنید.

\پاسخ{
با توجه به محدب بودن بزرگ‌ترین زیرمجموعۀ $k$ محسوس به صورت زیر عمل می‌کنیم. ابتدا $M$ را مرتب می‌کنیم. سپس تمام عناصر $M$ را به عنوان کوچکترین عنصر مجموعه امتحان می‌کنیم و بزرگ‌ترین عنصر را با جست‌وجوی دودویی به دست می‌آوریم. برای هر عنصر انتخاب شده به عنوان کوچکترین عنصر، $s$ را که حاصل‌جمع $k$ عنصر کوچکتر است را به دست می‌آوریم. اگر به ترتیب جلو برویم، $s$ را از مرحلۀ قبل با حذف یک عنصر و اضافه کردن یک عنصر در $O(1)$ به دست می‌آوریم. سپس به دنبال بزرگ‌ترین عنصری از $M$ می‌گردیم که کمتر مساوی $s$ باشد. تمام عناصر بین کوچکترین و بزرگترین عنصر انتخاب شده یک مجموعه $k$-محسوس است و بزرگترین این مجموعه‌ها، بزرگ‌ترین مجموعۀ $k$-محسوس $M$ خواهد بود.

زمان اجرای الگوریتم برای مرتب‌سازی $O(n\log n)$ و برای قسمت بعدی نیز $n\cdot (O(1) + O(\log n)) = O(n\log n)$ است. بنابراین زمان اجرای کل الگوریتم $O(n\log n)$ است.

نسخه سخت‌تری از سؤال که قرار بود مطرح شود به این صورت بود که اعداد $M$ دارای محدودیت هستند (اعداد صحیح بین ۱ تا $n^{100}$) و بنابراین زمان اجرایی مرتب‌سازی $O(n)$ می‌شود. همچنین با کمی تلاش می‌توان جست‌وجوی دودویی را حذف کرد و مکان‌نمایی برای بزرگ‌ترین عنصر $k$-محسوس درنظر گرفت و نشان داد این مکان‌نما در هر مرحله به عقب برنمی‌گردد. بنابراین قسمت دوم را نیز می‌توان در زمان $O(n)$ حل کرد.
	
}

\پایان{زیرمسائل}

%\شروع‌مسئله[12] 
%یک جدول $n\times n$ داریم که یک مار تمام خانه‌های آن را مشابه شکل زیر پوشانده است. ما (مثل مسیر هامیلتونی) تمام خانه‌ها را پوشانده است و از یک خانه دوبار عبور نکرده است. شما در هر مرحله می‌توانید به طور دلخواه یک زیرجدول را انتخاب کنید و این سؤال را بپرسید که تعداد بارهایی که مار مرز این جدول را رد می‌کند چند بار است؟ به طور مثال در شکل زیر ما چهار بار مرز زیرجدول را رد می‌کند و جواب سؤال زوج است. الگوریتمی طراحی کنید که با $O(n)$ پرسش سر و ته مار را پیدا کند.